\documentclass[../../main]{subfiles}
\begin{document}

This thesis investigated the feasibility of employing LLMs as decision components for solver selection in CP. The central research question concerned whether a general-purpose LLM, without domain-specific fine-tuning, can effectively map structural characteristics of combinatorial models to appropriate solver choices within a predefined portfolio.

The study focused on MiniZinc~\cite{MiniZinc} models derived from benchmark instances of the MiniZinc Challenge~\cite{MznProblems}. The experimental design systematically evaluated LLM-driven solver selection under multiple configurations, including different input representations (raw MiniZinc scripts, natural language descriptions, structured feature vectors, and natural-language abstractions of FlatZinc models), prompting strategies, and solver portfolio restrictions. Performance was assessed using standard ad-hoc metrics, with aim of enabling direct comparison with established single solvers.

The objective was to analyze the extent to which a general-purpose LLM can infer solver suitability from structural problem information. Particular attention was devoted to understanding the best problem representation in order to maximize the LLM's understanding of the input. Additionally, the study examined whether the LLM's internal reasoning, when made explicit, reflects a coherent understanding of solver categories and instance structure, and if it uses it accordingly in order to make proper choices.

Through this framework, the thesis aimed to provide a systematic and controlled assessment of LLM-based solver selection, identifying both its capabilities and its current limitations in the context of CP.

\section{Empirical findings}
As discussed in \Cref{chap:methodology} and detailed in \Cref{chap:expEval}, the experimental campaign was constrained by several practical limitations inherent to the use of commercial LLMs. In particular, API rate limits (cf. \Cref{tab:rate_limits_gemini}, \Cref{tab:rate_limits_groq}) imposed strict bounds on the number of admissible requests and on the maximum token throughput, thereby directly restricting both the scale and the granularity of the evaluation. These constraints became evident during the preliminary assessment over a broader set of LLMs (cf. \Cref{tab:combinedAllLLM}), where the number of testable instances remained limited. 

Rate limits were compounded by additional issues related to long-context inference. Large MiniZinc instances, especially those containing extensive arrays or parameter data, resulted in high token consumption and increased susceptibility to degradation phenomena such as context rot~\cite{ContextRot}. To mitigate these effects, input pre-processing strategies were introduced. In particular, script sanitization techniques were applied to reduce token usage, including systematic truncation of large arrays through controlled brute-force cutting, as described in \Cref{sec:scriptMan}. These measures aimed to preserve structural information while maintaining inputs within feasible context and rate boundaries.

Under this paradigm, initial results remained below the performance of the single best solver. However, the introduction of structured natural-language descriptions led to marginal improvements over na\"{\i}ve single-solver usage, as reported in \Cref{tab:combinedAllLLM}. Although the observed gains were limited in magnitude, they indicated that representational choices influence the effectiveness of LLM-based solver selection.

Motivated by these findings, more structured input transformations were explored. The feature extraction tool \texttt{mzn2feat}~\cite{mzn2feat} was integrated into the pipeline to generate fixed-dimensional representations encoding relevant characteristics of the MiniZinc models. These feature vectors were tested both in isolation and in combination with raw scripts or natural-language descriptions. Despite the increased regularity and dimensional stability of the input representation, no performance improvement beyond the previously obtained maximum score was observed (cf. \Cref{tab:combinedFeatVars}).

Subsequently, sampling temperature tuning was introduced to assess whether inference-time stochasticity influenced solver selection quality. The best-performing configurations were evaluated across a range of temperature values, as summarized in \Cref{tab:combinedFeatVarsTemp}. This adjustment yielded the highest overall score achieved in the unrestricted portfolio setting, indicating that controlled variability in model responses can affect selection outcomes.

Notwithstanding these incremental improvements, the overall performance gains remained limited. Moreover, a consistent selection pattern emerged: the LLM predominantly chose a small subset of globally dominant solvers across the majority of instances. While this strategy is not inherently ineffective, given the empirical dominance of those solvers, it reduces the discriminative value of the selection process and limits interpretability regarding instance-specific reasoning.

To further investigate the LLM's decision dynamics, the experimental setup was modified by restricting the solver portfolio. In this configuration, accurate and fine-grained discrimination among solvers was required to achieve competitive performance. Under these conditions, the limitations of the approach became more pronounced: the obtained scores were consistently below the single best solver within the restricted portfolio, as shown in \Cref{tab:CombinedSelSolversTop}. This outcome suggests that the model's selection strategy lacks sufficient sensitivity to intra-category solver differences when dominance effects are removed.

The observed influence of natural-language representations, together with the limitations identified in previous configurations, motivated the development of a dedicated transformation tool, \texttt{fzn2nl}~\cite{fzn2nl}, described in \Cref{chap:parser}. This tool generates a deterministic natural-language description of a problem instance starting from its FlatZinc representation. The transformation preserves structural information while eliminating most low-level syntactic artifacts, thereby providing a standardized and compact textual abstraction of the constraint model.

The output of \texttt{fzn2nl} was used as prompt input for solver selection over the restricted solver-set. Experimental results indicate that this representation did not yield performance improvements over the best previously obtained configurations. The achieved scores remained comparable to those of other input modalities (cf. \Cref{tab:combinedFznTemp}). However, the \texttt{fzn2nl} representation required the lowest token budget among all tested approaches, demonstrating improved efficiency in terms of context usage while maintaining competitive performance. 

Experiments conducted under the restricted solver portfolio further highlighted structural limitations in the selection strategy. Even in a setting where effective performance required diversified and instance-sensitive decisions, the LLM continued to concentrate its choices on a small subset of solvers, often selecting between two alternatives with high frequency. This behavior suggested that solver selection was not strongly driven by fine-grained instance characteristics.

To investigate the extent to which solver descriptions influence model decisions, additional controlled experiments were performed, as detailed in \Cref{chap:reasoning}. In a first setting, solver descriptions were permuted while preserving solver names. The results (cf. \Cref{tab:CombinedSwapped}) indicate that selections remained largely invariant under description swapping, suggesting that the model relied predominantly on solver names rather than on the associated technical descriptions.

To mitigate this bias, a subsequent experiment anonymized solver names, replacing them with neutral identifiers (e.g., a, b, c), while preserving their descriptions. Under this configuration, the LLM exhibited increased sensitivity to descriptive content, demonstrating a partial capacity to associate solver characteristics with problem structure. Nonetheless, the selection pattern remained highly concentrated on a limited subset of solvers, indicating persistent limitations in discriminative reasoning.

Finally, qualitative analysis was conducted by prompting the LLM to provide explicit justifications for its solver choices. The generated explanations reveal a coherent understanding of high-level problem characteristics, including objective structure, constraint types, and domain properties. The model consistently distinguishes among solver paradigms (e.g., CP, MILP, CP-SAT, \dots) and correctly associates broad structural traits with appropriate solver categories. However, the explanations also expose a limited capacity to accurately link specific instance features to empirically optimal solvers within a given category. Sensitivity to intra-category distinctions and relative solver competitiveness remains weak. Overall, the analysis suggests that, while the LLM demonstrates substantial structural comprehension, its selection strategy exhibits a systematic tendency toward ``authority-based'' decisions rather than evidence-driven differentiation.

\section{Technical contributions}
This thesis provided a set of methodological and technical contributions that extend beyond the specific experimental results and constitute reusable assets for subsequent research in solver selection.

First, a systematic experimental protocol for evaluating LLM-based solver selection under controlled conditions has been defined. The protocol specifies standardized input representations, solver portfolios, prompting templates, temperature configurations, with evaluation metrics based on previous researches~\cite{evaluationMetaSolvers}. It further incorporates controlled comparisons against single solvers, enabling the isolation of dominance effects and the assessment of genuine discriminative capability. This framework ensures reproducibility and facilitates rigorous benchmarking of alternative LLM configurations or future models.

Second, a structured comparative analysis of prompting strategies has been conducted. The study evaluates raw MiniZinc scripts, natural-language reformulations, structured feature vectors, and hybrid combinations thereof. In addition, sampling temperature, have been systematically explored. This analysis provides empirical evidence regarding the impact of representational and stochastic factors on solver selection performance, establishing a reference baseline for future investigations into prompt engineering and representation learning in the context of CP.

Third, this thesis introduces \texttt{fzn2nl}~\cite{fzn2nl}, a tool for deterministic natural-language extraction from FlatZinc models. Its development required modifications to the compilation pipeline and the design of a rule-based transformation layer capable of converting low-level FlatZinc representations into structured textual descriptions while preserving semantic content. The resulting representation is compact, token-efficient, and model-agnostic. As a preprocessing component, \texttt{fzn2nl} can be integrated into other LLM-based workflows that require interpretable abstractions of MiniZinc and FlatZinc programs.

Finally, an interpretative framework for analyzing LLM reasoning in solver-selection tasks has been established. Through controlled perturbation experiments, such as description permutation and solver anonymization, and qualitative examination of generated explanations, this approach provides a methodological basis for evaluating not only predictive performance but also the internal consistency and structural grounding of LLM decisions in this context.

\section{Limitations and Future Work}
The experimental findings and methodological framework presented in this thesis must be interpreted in light of several structural constraints that limit their generalizability.

As discussed throughout the manuscript, the experimental campaign was significantly constrained by interaction budgets imposed by API rate limits and token throughput restrictions. Even in the case of \texttt{GPT-5.2}, it was not possible to conduct exhaustive evaluations across all instance sets, temperature configurations, or repeated trials. These operational constraints reduced the statistical depth of certain analyses and limited the systematic exploration of performance variability across multiple runs.

An additional limitation concerns the intrinsic non-determinism of LLMs. Solver selection outcomes may differ across identical prompt inputs due to stochastic inference mechanisms. This variability affects the stability of individual predictions and, consequently, constrains the generality of the evaluation.

The evaluation domain itself introduces further boundaries. All experiments were conducted on benchmark instances derived from the MiniZinc Challenge, which constitute a standardized and diverse testbed for CP. However, this benchmark suite does not encompass the full spectrum of modeling paradigms or real-world application domains. The conclusions therefore primarily characterize LLM-based solver selection within the MiniZinc ecosystem.

These limitations suggest several directions for future research. Extending the experimental framework to larger and more heterogeneous benchmark collections, and evaluating more advanced LLM architectures with increased reasoning capacity and extended context windows, would enable a more robust assessment of scalability and generalization. Broader empirical coverage would allow a more precise identification of the conditions under which LLM-assisted solver selection may achieve competitive performance.

It should also be emphasized that the objective of this research was to investigate the potential of general-purpose LLMs in the context of solver selection. The absence of task-specific adaptation highlights the opportunity to explore fine-tuning strategies based on solver-performance data. Such approaches could enable models to internalize empirical competitiveness and reduce reliance on implicit priors acquired during pretraining, thereby strengthening the association between structural instance features and solver choice.

Overall, while current general-purpose LLMs may not constitute competitive alternatives in the field of solver selection, they exhibit non-trivial structural understanding of constraint models and solver paradigms. The results delineate both the present limitations and a plausible trajectory for the integration of LLM-based components into CP workflows, contingent upon further methodological refinement and empirical validation.

\end{document}